\section{第三层：数学与解释结构}

\subsection{从局部更新到整体机制}

涌现语言的数学价值，不在于为日常生活制造精确的闭式公式。
它的价值在于澄清解释的架构。

\begin{proposition}[微观到宏观的反馈]
如果重复的局部更新会修改未来转变的成本，那么由此形成的宏观模式就不只是描述性的。
它在模型中会变得具有因果相关性，因为它改变了后续轨迹何者廉价、何者稳定、何者可达。
\end{proposition}

\begin{discussion}
这正是因果涌现语言背后的解释点。
一旦一个成熟机制改变了后续转变几何，这个机制就值得被视为不仅仅是一个摘要标签。
它已经成为一个与控制相关的描述层级。
\end{discussion}

\subsection{机制、盆地与重新进入成本}

本章并不要求沉重的方程，但它确实要求一个更严格的数学图景。
设一个行为机制可被表示为状态空间中的一个区域，并具有以下性质：
\begin{enumerate}[nosep]
  \item \textbf{进入成本：} 进入该区域所需的协调量有多大；
  \item \textbf{保持性：} 一旦进入后，轨迹留在其中的可能性有多高；
  \item \textbf{重新进入经济性：} 中断之后，系统能以多低成本回到该区域；
  \item \textbf{竞争性抑制：} 已建立机制会在多大程度上削弱竞争路径的吸引力。
\end{enumerate}

这样一来，框架形成就可以被解释为一个逐步完成以下变化的过程：
\begin{itemize}[nosep]
  \item 降低某一路径的进入成本；
  \item 加深或锐化其吸引盆；
  \item 通过反复协调使用储存路径记忆；
  \item 提高竞争路径的干扰成本。
\end{itemize}

\subsection{为什么涌现不只是聚合}

如果只是简单聚合，就意味着许多次重复只是机械相加。
而本章中的涌现比这更强。
重复行为并不只是累积；它们会重构未来的问题本身。

\begin{lemma}[重构判据]
如果重复的局部事件改变了未来转变景观本身，那么由此形成的高层模式所做的事情就不只是对这些事件进行汇总。
\end{lemma}

\begin{proof}
设想局部事件只是相加，而不重构后续选择。
那么未来决策问题在结构上并未改变，重复事件也就没有创造出新的解释层级。
但如果这些事件降低了某一路径的重新进入成本、增强了其保持性，或提高了针对竞争路径的干扰强度，
那么未来问题就已经不同于没有这段历史时的情形。
因此，这个宏观模式已经成为模型相关的实体。
\end{proof}

\subsection{数学升级路径}

本章刻意比第~\ref{ch:advanced-dynamics} 章更轻一些，
但它以三种方式为该章做准备：
\begin{enumerate}[nosep]
  \item 亚稳态说明了为什么有用框架必须能够持续，而又不至于不可改变；
  \item 吸引子语言说明了重复使用如何把某一路径变成默认盆地；
  \item 算子与控制视角将在后文帮助描述：干预所作用的不是一张白纸，而是一个对历史敏感的系统。
\end{enumerate}